\documentclass[tikz,border=8pt]{standalone}
\usepackage{tikz}
\usepackage{amsmath}
\usepackage{amssymb}
\usepackage{pgfplots}
\pgfplotsset{compat=1.18}
\usepackage{ctex}
\usetikzlibrary{calc, positioning, arrows.meta, decorations.pathreplacing}

\begin{document}
\begin{tikzpicture}[>=Stealth]

%% ===== 主图：平行四边形 =====
\begin{scope}[xshift=0cm]

  % 背景网格
  \draw[gray!25, thin] (-0.3,-0.3) grid (4.3,3.3);
  % 坐标轴
  \draw[->, gray!60, thin] (-0.3,0) -- (4.7,0) node[right, font=\scriptsize] {$x$};
  \draw[->, gray!60, thin] (0,-0.3) -- (0,3.6) node[above, font=\scriptsize] {$y$};

  % 平行四边形顶点：O=(0,0), A=a1=(3,1), B=a2=(1,2), C=a1+a2=(4,3)
  \fill[blue!18, opacity=0.7] (0,0) -- (3,1) -- (4,3) -- (1,2) -- cycle;
  \draw[blue!50, thin] (3,1) -- (4,3);
  \draw[blue!50, thin] (1,2) -- (4,3);

  % 列向量 a1
  \draw[->, thick, blue!80!black] (0,0) -- (3,1)
    node[below right, font=\small] {$\mathbf{a}_1=(3,1)$};
  % 列向量 a2
  \draw[->, thick, teal!80!black] (0,0) -- (1,2)
    node[left, font=\small] {$\mathbf{a}_2=(1,2)$};

  % 原点标注
  \fill (0,0) circle (1.8pt) node[below left, font=\scriptsize] {$O$};

  % 面积标注（中心）
  \node[font=\normalsize\bfseries, blue!70!black] at (2.0, 1.5)
    {面积 $= |\det A| = 5$};

  % 子图标题
  \node[font=\small\bfseries] at (2.0, -0.85) {$2\times2$ 行列式 $=$ 平行四边形面积};

\end{scope}

%% ===== 公式框 =====
\node[font=\small, fill=white, draw=gray!50, thin, rounded corners=2pt,
      inner sep=7pt, align=center] at (2.0, -1.9) {
  $\det\begin{pmatrix}3 & 1\\1 & 2\end{pmatrix} = 3\cdot 2 - 1\cdot 1 = 5$
};

%% ===== 右侧小图：行列式正负与方向 =====
\begin{scope}[xshift=6.8cm, yshift=0.2cm]

  % 小图标题
  \node[font=\small\bfseries] at (1.5, 3.3) {行列式符号与方向};

  %% 上半：正（逆时针）
  \begin{scope}[yshift=1.6cm]
    \draw[gray!25, thin] (-0.2,-0.1) grid (2.7,1.6);
    \draw[->, gray!60, thin] (-0.2,0) -- (2.9,0) node[right, font=\tiny] {$x$};
    \draw[->, gray!60, thin] (0,-0.1) -- (0,1.7) node[above, font=\tiny] {$y$};

    % 正向：a1=(2,0.3), a2=(0.3,1.2) → 逆时针
    \fill[blue!15, opacity=0.7] (0,0) -- (2,0.3) -- (2.3,1.5) -- (0.3,1.2) -- cycle;
    \draw[->, thick, blue!70!black] (0,0) -- (2,0.3)
      node[below, font=\tiny] {$\mathbf{a}_1$};
    \draw[->, thick, teal!70!black] (0,0) -- (0.3,1.2)
      node[left, font=\tiny] {$\mathbf{a}_2$};

    % 逆时针弧线
    \draw[->, thick, orange!80!black] (0.55,0.1) arc[start angle=10, end angle=68, radius=0.5cm];
    \node[font=\tiny, orange!80!black] at (0.95,0.45) {逆时针};
    \node[font=\tiny\bfseries, blue!70!black] at (1.5,-0.38) {$\det > 0$};
  \end{scope}

  %% 下半：负（顺时针）
  \begin{scope}[yshift=-0.4cm]
    \draw[gray!25, thin] (-0.2,-0.1) grid (2.7,1.6);
    \draw[->, gray!60, thin] (-0.2,0) -- (2.9,0) node[right, font=\tiny] {$x$};
    \draw[->, gray!60, thin] (0,-0.1) -- (0,1.7) node[above, font=\tiny] {$y$};

    % 负向：交换 a1 和 a2 → 顺时针
    \fill[red!15, opacity=0.7] (0,0) -- (0.3,1.2) -- (2.3,1.5) -- (2,0.3) -- cycle;
    \draw[->, thick, red!70!black] (0,0) -- (0.3,1.2)
      node[left, font=\tiny] {$\mathbf{a}_1'$};
    \draw[->, thick, orange!50!black] (0,0) -- (2,0.3)
      node[below, font=\tiny] {$\mathbf{a}_2'$};

    % 顺时针弧线
    \draw[->, thick, red!80!black] (0.55,0.1) arc[start angle=68, end angle=10, radius=0.5cm];
    \node[font=\tiny, red!80!black] at (1.05,0.45) {顺时针};
    \node[font=\tiny\bfseries, red!70!black] at (1.5,-0.38) {$\det < 0$};
  \end{scope}

\end{scope}

%% 整体标题
\node[font=\large\bfseries] at (4.0, 4.2)
  {行列式的几何意义：面积与方向};

\end{tikzpicture}
\end{document}
